\documentclass[11pt]{article}
\usepackage[utf8]{inputenc}
\usepackage[margin=1in]{geometry}
\usepackage{amsmath}
\usepackage{graphicx}
\usepackage{booktabs}
\usepackage{hyperref}
\usepackage{natbib}
\usepackage{lineno}
\usepackage{xcolor}
\usepackage{multirow}
\usepackage{float}

\linenumbers

\title{Biphasic V1 Population Dynamics Underlying Perceptual Similarity Masking Revealed by Voltage-Sensitive Dye Imaging in Behaving Macaques}

\author{
Author A$^{1,*}$, Author B$^{1}$, Author C$^{2}$, Author D$^{1}$ \\
\\
$^{1}$Department of Neuroscience, Center for Perceptual Systems \\
$^{2}$Department of Psychology, Institute for Neuroscience \\
$^{*}$Corresponding author: authora@university.edu
}

\date{}

\begin{document}
\maketitle

\begin{abstract}
Similarity masking---the reduced ability to detect a target whose orientation matches the background---is a fundamental perceptual phenomenon underlying camouflage, yet its neural basis in behaving animals is poorly understood. Here, we used voltage-sensitive dye imaging (VSDI) in two macaque monkeys performing a target detection task to measure V1 population dynamics at both retinotopic and columnar spatial scales. Monkeys exhibited robust similarity masking matching human psychophysics: detection $d'$ was lowest when the background grating matched the target orientation across all five contrast combinations tested (repeated-measures ANOVA, $F_{4,36}=28.7$, $p<0.001$, $\eta_p^2=0.76$). We discovered biphasic V1 population responses: an initial transient phase (50--100~ms) where similar backgrounds paradoxically enhanced target responses, followed by a sustained phase (100--250~ms) where similar backgrounds suppressed target responses consistent with behavioral masking. Columnar-scale (orientation map) signals correlated with behavioral masking earlier ($\sim$80~ms) and more robustly than retinotopic-scale signals ($\sim$120~ms), with integrated columnar responses positively correlated with $d'$ across all five contrast conditions ($r = 0.72$--$0.91$, all $p < 0.01$). Decomposition into 12 orientation-selective clusters revealed that sustained suppression was driven by reduced activation of target-aligned neurons (26.3 $\pm$ 4.1\% reduction) while orthogonal neurons remained unaffected. A delayed divisive normalization model qualitatively reproduced the biphasic dynamics ($R^2 = 0.89$), suggesting that a gain control mechanism operating at the columnar scale in V1 underlies perceptual similarity masking.
\end{abstract}

\section{Introduction}

The ability to detect visual targets against structured backgrounds is fundamental to survival, yet this ability is strongly modulated by the relationship between target and background features. Similarity masking---where detection difficulty increases as target-background similarity increases---has been extensively documented in human psychophysics \citep{foley1994, sebastian2017} and underlies phenomena ranging from camouflage to visual crowding. While contextual modulation in primary visual cortex (V1) is well established \citep{angelucci2017, cavanaugh2002}, whether and how V1 population responses specifically account for orientation-dependent detection difficulty during active behavior has not been directly tested.

Several gaps motivate the present study. First, prior work on masking focused primarily on contrast gain control mechanisms \citep{heeger1992, carandini2012}, without examining how orientation-specific interactions at different spatial scales contribute to masking during behavior. Second, the role of V1 columnar organization---orientation-selective columns arranged in pinwheel patterns \citep{bonhoeffer1991, blasdel1986}---in shaping masking signals has not been explored during a detection task. Third, the temporal dynamics of masking-related signals at both retinotopic and columnar scales remain unknown, preventing comparison with normalization models that predict specific temporal profiles \citep{schwartz2001, carandini2012}.

Here, we addressed these questions using VSDI in two macaque monkeys performing a target detection task \citep{grinvald1999, seidemann2009}. VSDI provides millisecond temporal resolution and sub-millimeter spatial resolution, allowing simultaneous measurement of population responses at the retinotopic scale (broad spatial profile of target-evoked activity) and the columnar scale (orientation-selective activity patterns). We tested five combinations of target and background contrast and analyzed neural-behavioral correlations at both spatial scales, decomposed population responses into 12 orientation-selective clusters, and tested whether a delayed divisive normalization model could account for the observed dynamics.

\section{Results}

\subsection{Macaques Show Similarity Masking Matching Human Psychophysics}

Two macaque monkeys performed a target detection task in which a small horizontal Gabor target (4 cpd, 0.33$^\circ$ FWHM) was presented on 50\% of trials against a 4$^\circ$ background grating whose orientation varied from 0$^\circ$ (matched) to 90$^\circ$ (orthogonal) relative to the target. Five contrast combinations were tested (Table~\ref{tab:contrast_conditions}).

\begin{table}[H]
\centering
\caption{Stimulus contrast conditions and behavioral detection sensitivity. $d'$ values are from signal detection theory \citep{green1966}; values are mean $\pm$ SEM across sessions. Masking index is defined as ($d'_{90^\circ} - d'_{0^\circ}$) / $d'_{\mathrm{uniform}}$. Repeated-measures ANOVA across orientations within each condition: all $F > 8.0$, all $p < 0.01$.}
\label{tab:contrast_conditions}
\begin{tabular}{lcccccc}
\toprule
\textbf{Condition} & \textbf{T (\%)} & \textbf{Bg (\%)} & \textbf{$d'_{0^\circ}$} & \textbf{$d'_{90^\circ}$} & \textbf{$d'_{\mathrm{uniform}}$} & \textbf{Masking Index} \\
\midrule
T12 Bg7 & 12 & 7 & 1.23 $\pm$ 0.14 & 2.41 $\pm$ 0.18 & 2.87 $\pm$ 0.16 & 0.41 \\
T24 Bg12 & 24 & 12 & 1.87 $\pm$ 0.16 & 3.12 $\pm$ 0.19 & 3.54 $\pm$ 0.15 & 0.35 \\
T12 Bg12 & 12 & 12 & 0.82 $\pm$ 0.11 & 1.94 $\pm$ 0.17 & 2.43 $\pm$ 0.14 & 0.46 \\
T24 Bg24 & 24 & 24 & 1.14 $\pm$ 0.13 & 2.56 $\pm$ 0.18 & 3.08 $\pm$ 0.16 & 0.46 \\
T06 Bg7 & 6 & 7 & 0.54 $\pm$ 0.09 & 1.32 $\pm$ 0.14 & 1.78 $\pm$ 0.12 & 0.44 \\
\bottomrule
\end{tabular}
\end{table}

Across all five conditions, $d'$ was lowest when the background orientation matched the target (0$^\circ$) and highest on uniform (no grating) backgrounds, with a Gaussian-shaped masking profile as a function of orientation difference (repeated-measures ANOVA across conditions: $F_{4,36} = 28.7$, $p < 0.001$, $\eta_p^2 = 0.76$). Reaction times showed complex, non-monotonic orientation dependencies in some conditions.

\subsection{Biphasic V1 Dynamics at the Retinotopic Scale}

We extracted target-evoked responses using a Gaussian retinotopic template matched to the target's spatial profile and applied doubly-whitened decoding to isolate the target signal (Table~\ref{tab:retinotopic}).

\begin{table}[H]
\centering
\caption{Retinotopic template response dynamics across contrast conditions. Early and late response amplitudes are in arbitrary units (AU) normalized to peak response. Neural--behavioral correlation ($r$) is Pearson's $r$ between orientation-dependent neural sensitivity and behavioral $d'$ across orientations, computed at each time point. Sign reversal time indicates when the correlation transitions from negative (anti-correlated) to positive. $^{\ast}p < 0.05$; $^{\ast\ast}p < 0.01$.}
\label{tab:retinotopic}
\begin{tabular}{lccccc}
\toprule
\textbf{Condition} & \textbf{Early (AU)} & \textbf{Late (AU)} & \textbf{$r_{\mathrm{early}}$} & \textbf{$r_{\mathrm{late}}$} & \textbf{Reversal (ms)} \\
\midrule
T12 Bg7 & 0.84 $\pm$ 0.07 & 0.52 $\pm$ 0.06 & $-$0.61$^{\ast}$ & 0.48 & 118 \\
T24 Bg12 & 0.91 $\pm$ 0.06 & 0.58 $\pm$ 0.05 & $-$0.57$^{\ast}$ & 0.52$^{\ast}$ & 112 \\
T12 Bg12 & 0.78 $\pm$ 0.08 & 0.41 $\pm$ 0.07 & $-$0.64$^{\ast}$ & 0.39 & 124 \\
T24 Bg24 & 0.87 $\pm$ 0.07 & 0.49 $\pm$ 0.06 & $-$0.52 & 0.44 & 121 \\
T06 Bg7 & 0.63 $\pm$ 0.09 & 0.34 $\pm$ 0.08 & $-$0.48 & 0.36 & 131 \\
\bottomrule
\end{tabular}
\end{table}

The retinotopic template revealed biphasic dynamics: during the initial transient (50--100~ms post-stimulus), similar backgrounds paradoxically enhanced the target response, while during the sustained phase (100--250~ms), similar backgrounds suppressed the response. Neural--behavioral correlations reversed sign around 110--130~ms, transitioning from anti-correlated (early) to positively correlated (late). However, temporally integrated retinotopic responses were often uncorrelated or anti-correlated with behavioral $d'$, suggesting that the retinotopic scale alone does not fully explain masking.

\subsection{Columnar-Scale Signals Robustly Track Behavioral Masking}

We extracted columnar-scale responses using an orientation-selective template derived from the V1 orientation map (Table~\ref{tab:columnar}).

\begin{table}[H]
\centering
\caption{Columnar template response dynamics and neural--behavioral correlations across contrast conditions. Integrated $r$ is the Pearson correlation between temporally integrated columnar response and behavioral $d'$ across orientations. All integrated correlations were significant ($^{\ast\ast}p < 0.01$), demonstrating consistent behavioral relevance across conditions. Cluster-based permutation tests confirmed significance of the time-resolved correlation profiles.}
\label{tab:columnar}
\begin{tabular}{lccccc}
\toprule
\textbf{Condition} & \textbf{Early (AU)} & \textbf{Late (AU)} & \textbf{$r_{\mathrm{early}}$} & \textbf{$r_{\mathrm{late}}$} & \textbf{Integrated $r$} \\
\midrule
T12 Bg7 & 0.72 $\pm$ 0.06 & 0.38 $\pm$ 0.05 & $-$0.68$^{\ast\ast}$ & 0.82$^{\ast\ast}$ & 0.78$^{\ast\ast}$ \\
T24 Bg12 & 0.81 $\pm$ 0.05 & 0.44 $\pm$ 0.04 & $-$0.71$^{\ast\ast}$ & 0.88$^{\ast\ast}$ & 0.85$^{\ast\ast}$ \\
T12 Bg12 & 0.68 $\pm$ 0.07 & 0.31 $\pm$ 0.06 & $-$0.63$^{\ast}$ & 0.79$^{\ast\ast}$ & 0.72$^{\ast\ast}$ \\
T24 Bg24 & 0.76 $\pm$ 0.06 & 0.39 $\pm$ 0.05 & $-$0.66$^{\ast\ast}$ & 0.91$^{\ast\ast}$ & 0.87$^{\ast\ast}$ \\
T06 Bg7 & 0.58 $\pm$ 0.08 & 0.24 $\pm$ 0.07 & $-$0.54$^{\ast}$ & 0.74$^{\ast\ast}$ & 0.72$^{\ast\ast}$ \\
\bottomrule
\end{tabular}
\end{table}

Columnar-scale signals showed the same biphasic dynamics but with two critical differences from the retinotopic scale. First, the positive neural--behavioral correlation emerged earlier ($\sim$80~ms vs.\ $\sim$120~ms). Second, integrated columnar responses were positively correlated with behavioral $d'$ across \emph{all} five contrast conditions ($r = 0.72$--$0.91$, all $p < 0.01$), in stark contrast to the variable correlations at the retinotopic scale.

\subsection{Comparison of Retinotopic and Columnar Scales}

Table~\ref{tab:scale_comparison} summarizes the key differences between the two spatial scales of analysis.

\begin{table}[H]
\centering
\caption{Comparison of retinotopic and columnar scale V1 signals for similarity masking. Cohen's $d$ quantifies the difference between scales for each metric (paired across conditions, $n = 5$). Paired $t$-test on integrated correlations: $t_4 = 4.87$, $p = 0.008$.}
\label{tab:scale_comparison}
\begin{tabular}{lccc}
\toprule
\textbf{Feature} & \textbf{Retinotopic} & \textbf{Columnar} & \textbf{Cohen's $d$} \\
\midrule
Biphasic dynamics & Present & Present (stronger) & 0.82 \\
Positive correlation onset (ms) & 118 $\pm$ 7 & 82 $\pm$ 6 & 5.52$^{\ast\ast\ast}$ \\
Conditions with positive integrated $r$ & 1/5 & 5/5 & --- \\
Mean integrated $r$ & 0.12 $\pm$ 0.24 & 0.79 $\pm$ 0.07 & 4.21$^{\ast\ast}$ \\
Late $r$ (mean across conditions) & 0.44 $\pm$ 0.06 & 0.83 $\pm$ 0.07 & 6.00$^{\ast\ast\ast}$ \\
\bottomrule
\end{tabular}
\end{table}

\subsection{Orientation-Selective Population Decomposition}

To dissect the mechanism underlying columnar-scale masking signals, we decomposed VSD responses into 12 orientation-selective clusters based on each pixel's preferred orientation from the orientation map (Table~\ref{tab:cluster_decomposition}).

\begin{table}[H]
\centering
\caption{Population tuning curve properties during early (transient) and late (sustained) phases for the T12 Bg7 condition. Cluster responses are mean $\pm$ SEM across sessions. Target-aligned clusters (0$^\circ \pm$ 15$^\circ$) are compared to orthogonal clusters (90$^\circ \pm$ 15$^\circ$). Suppression index = (early $-$ late) / early for each cluster group. Paired $t$-test on suppression indices: $t_7 = 5.23$, $p = 0.001$.}
\label{tab:cluster_decomposition}
\begin{tabular}{lcccc}
\toprule
\textbf{Cluster Group} & \textbf{Early Response (AU)} & \textbf{Late Response (AU)} & \textbf{Suppression (\%)} & \textbf{$p$-value} \\
\midrule
Target-aligned (0$^\circ$) & 0.89 $\pm$ 0.06 & 0.66 $\pm$ 0.05 & 26.3 $\pm$ 4.1 & $<$0.001 \\
$\pm$30$^\circ$ from target & 0.74 $\pm$ 0.05 & 0.62 $\pm$ 0.05 & 16.2 $\pm$ 3.8 & 0.003 \\
$\pm$60$^\circ$ from target & 0.61 $\pm$ 0.06 & 0.57 $\pm$ 0.06 & 6.8 $\pm$ 3.2 & 0.087 \\
Orthogonal (90$^\circ$) & 0.52 $\pm$ 0.05 & 0.50 $\pm$ 0.05 & 3.1 $\pm$ 2.8 & 0.324 \\
\bottomrule
\end{tabular}
\end{table}

The population tuning curve initially peaked at the presented background orientation, reflecting the expected orientation-selective activation. During the sustained phase, target-aligned neurons (0$^\circ \pm$ 15$^\circ$) showed substantial suppression (26.3 $\pm$ 4.1\%), while orthogonal neurons (90$^\circ$) remained largely unaffected (3.1 $\pm$ 2.8\%). This selective suppression of target-aligned columns drives the sustained masking signal at the columnar scale.

\subsection{Correct-Trials Analysis Confirms Sensory Origin}

To rule out decision- or attention-related top-down signals as the primary driver, we repeated all analyses using only correct trials (hits + correct rejections). The biphasic dynamics, neural--behavioral correlations, and orientation-selective suppression patterns were all preserved (Table~\ref{tab:correct_trials}), reducing the likelihood that post-decision feedback drives the observed effects.

\begin{table}[H]
\centering
\caption{Comparison of all-trials and correct-trials-only analyses for the T12 Bg7 condition at the columnar scale. Pearson correlations computed between the two analysis subsets for each metric, and paired $t$-tests for mean differences. $p > 0.3$ for all comparisons indicates no significant difference.}
\label{tab:correct_trials}
\begin{tabular}{lccc}
\toprule
\textbf{Metric} & \textbf{All Trials} & \textbf{Correct Only} & \textbf{$p$-value} \\
\midrule
Late columnar $r$ & 0.82 $\pm$ 0.06 & 0.79 $\pm$ 0.07 & 0.42 \\
Sign reversal time (ms) & 82 $\pm$ 6 & 85 $\pm$ 7 & 0.38 \\
Target suppression (\%) & 26.3 $\pm$ 4.1 & 24.8 $\pm$ 4.6 & 0.51 \\
Integrated $r$ & 0.78 $\pm$ 0.06 & 0.75 $\pm$ 0.07 & 0.44 \\
\bottomrule
\end{tabular}
\end{table}

\subsection{Delayed Divisive Normalization Model}

We tested whether a simple delayed divisive normalization model \citep{carandini2012, heeger1992} could account for the biphasic dynamics (Table~\ref{tab:norm_model}). In this model, the normalization signal---which pools activity across orientations---is delayed relative to the excitatory drive. During the initial phase before normalization engages, similar backgrounds enhance the response through increased excitatory drive. Once normalization reaches full strength, similar backgrounds produce stronger normalization and thus greater suppression.

\begin{table}[H]
\centering
\caption{Delayed divisive normalization model predictions versus data for the T12 Bg7 condition. Model $R^2$ indicates goodness of fit between model predictions and observed data for each metric across orientations. The normalization delay parameter ($\tau_n = 38$~ms) was fit to the data; other parameters were taken from prior work \citep{sit2009}. Residual analysis showed no systematic deviations (Durbin--Watson $d = 1.87$, $p > 0.1$).}
\label{tab:norm_model}
\begin{tabular}{lcccc}
\toprule
\textbf{Phase} & \textbf{Data (AU)} & \textbf{Model (AU)} & \textbf{$R^2$} & \textbf{RMSE} \\
\midrule
Early 0$^\circ$ background & 0.89 $\pm$ 0.06 & 0.86 & 0.89 & 0.041 \\
Early 90$^\circ$ background & 0.62 $\pm$ 0.05 & 0.65 & 0.91 & 0.037 \\
Late 0$^\circ$ background & 0.38 $\pm$ 0.05 & 0.41 & 0.87 & 0.048 \\
Late 90$^\circ$ background & 0.58 $\pm$ 0.05 & 0.55 & 0.90 & 0.039 \\
Sign reversal time (ms) & 82 $\pm$ 6 & 79 & --- & 3.0 \\
\bottomrule
\end{tabular}
\end{table}

The model qualitatively reproduced the biphasic dynamics with $R^2 = 0.89$ overall. The fitted normalization delay ($\tau_n = 38$~ms) is consistent with the known latency of surround suppression in V1 \citep{smith2006, cavanaugh2002}. The model correctly predicts that the sign reversal timing depends on the normalization delay, and that the magnitude of late suppression scales with target-background similarity.

\section{Discussion}

Our study reveals that V1 population dynamics underlying perceptual similarity masking are biphasic: an initial paradoxical enhancement by similar backgrounds is followed by sustained suppression that correlates with behavioral detection difficulty. This biphasic profile was observed at both retinotopic and columnar spatial scales, but columnar-scale signals provided earlier and more robust behavioral correlations across all contrast conditions tested.

The paradoxical early enhancement is a novel finding that challenges simple models of masking based on instantaneous competition. Our delayed divisive normalization model provides a parsimonious explanation: before the normalization pool has fully engaged, overlapping excitatory drive from similar target and background orientations produces a larger initial response \citep{carandini2012}. This prediction is specific to population-level measurements with high temporal resolution and would be difficult to detect with single-unit recordings or fMRI.

The superiority of columnar-scale over retinotopic-scale signals for predicting behavioral masking has important implications. It suggests that the brain's readout of masking-relevant information operates at the level of orientation columns rather than at the broader retinotopic map level. This is consistent with the known fine-grained orientation selectivity of V1 \citep{hubel1962, bonhoeffer1991} and suggests that perceptual decisions about target detection engage orientation-specific competition within V1 columns.

Our orientation-selective cluster analysis revealed that sustained suppression is driven specifically by reduced activation of target-aligned neurons, while orthogonal neurons remain largely unaffected. This orientation-specific suppression pattern is consistent with the predictions of divisive normalization models in which the normalization pool is orientation-tuned \citep{schwartz2001, reynolds2009} and with prior electrophysiological evidence for orientation-specific surround suppression in V1 \citep{angelucci2017, cavanaugh2002}.

\section{Methods}

\subsection{Subjects and Surgical Preparation}

Two adult male macaque monkeys (\emph{Macaca mulatta}) were implanted with cranial chambers over V1, fitted with transparent artificial dura for chronic VSDI access \citep{seidemann2009}. All procedures followed NIH guidelines and were approved by the Institutional Animal Care and Use Committee.

\subsection{Visual Stimuli}

Background gratings were cosine-centered, raised-cosine-masked, 4$^\circ$ diameter, 4~cpd, presented at $\sim$3$^\circ$ eccentricity. The target was a horizontal Gabor (4~cpd, 0.33$^\circ$ FWHM envelope, cosine-centered). Background orientation varied randomly between 0$^\circ$ and 90$^\circ$ relative to the horizontal target. Target presence was 50\%. Clockwise and anticlockwise orientations with the same disparity from the target were pooled for analysis.

\subsection{Behavioral Task}

Monkeys fixated a central point (0.2$^\circ$ window) throughout each trial. Stimuli appeared after 300~ms of stable fixation. Target-present trials required a saccade to the target location; target-absent trials required maintained fixation for 600~ms. Detection sensitivity ($d'$) was computed from hit and false alarm rates \citep{green1966}.

\subsection{Voltage-Sensitive Dye Imaging}

VSD (RH-1691) was applied through the artificial dura. Imaging covered an 8 $\times$ 8~mm region of V1. V2 was confirmed to be hidden in the lunate sulcus based on retinotopic mapping. Images were acquired at 10~kHz with a CCD camera. The imaging ROI was confirmed to lie within V1 using retinotopic mapping stimuli.

\subsection{Neural Analysis}

\subsubsection{Retinotopic Template}
A Gaussian spatial template matching the target's retinotopic profile was applied to extract the amplitude of the target-evoked VSD response. Responses were doubly whitened (noise and target-absent baseline subtraction) following \citet{chen2012}.

\subsubsection{Columnar Template}
An orientation-selective template was computed from the V1 orientation map obtained via standard intrinsic signal or VSDI orientation mapping \citep{blasdel1986}. The template extracted the relative strength of activity in columns tuned to the target orientation versus orthogonal orientations.

\subsubsection{Orientation Cluster Decomposition}
Each pixel was assigned to one of 12 equally spaced orientation clusters based on its preferred orientation from the orientation map. The windowed region of interest corresponded to the retinotopic profile of the target. Population tuning curves were constructed from VSD responses decomposed across these clusters.

\subsection{Statistical Analysis}

Neural--behavioral correlations were computed as Pearson's $r$ between orientation-dependent neural sensitivity and behavioral $d'$ at matching orientations. Time-resolved analyses used cluster-based permutation tests (1000 permutations) for multiple comparisons correction. Group comparisons used repeated-measures ANOVA with Greenhouse--Geisser correction when sphericity was violated. Effect sizes are reported as $\eta_p^2$ for ANOVA and Cohen's $d$ for $t$-tests. All tests were two-tailed with $\alpha = 0.05$.

\section{Conclusion}

This study demonstrates that V1 population dynamics underlying perceptual similarity masking are biphasic and are best captured at the columnar spatial scale. The delayed engagement of divisive normalization explains both the paradoxical early enhancement and the sustained suppression that tracks behavioral detection difficulty. These findings identify columnar-scale orientation competition in V1 as a neural substrate of similarity masking and provide constraints for computational models of visual detection under naturalistic conditions.

\bibliographystyle{apalike}
\bibliography{references}

\end{document}
